\section{Conclusions}

This randomized, grader-blinded trial found no statistically significant or clinically meaningful difference in post-operative periorbital bruising between treatment conditions across all outcomes and all statistical models.

For bruising depth, the estimated condition effect was $\beta = -0.034$ (LMM), OR $= 0.822$ (CLMM), and $\beta = -0.034$ (GEE-Ordinal), all with $p > 0.70$. For bruising extent, the estimates were $\beta = 0.029$ (LMM), OR $= 1.006$ (CLMM), and $\beta = 0.029$ (GEE-Ordinal), all with $p > 0.82$. The composite outcome (depth + extent) showed $\beta = -0.005$ ($p = 0.990$). Nonparametric analyses (Mann--Whitney, permutation tests) and multiple-comparison corrections confirmed the null finding. All Cohen's $d$ values were below 0.12, classified as negligible by standard conventions, and Cliff's $\delta$ values were below 0.09.

Notwithstanding this consistency, the trial was severely underpowered. Post-hoc power ranged from 5.5\% (extent) to 7.2\% (depth), with minimum detectable effects of approximately 1.0~SD. These values indicate that the study could only have detected implausibly large treatment effects and that the null result does not constitute strong evidence against the intervention. Sample sizes of 571 to 2{,}786 patients per group would be required to detect effects of the observed magnitude at 80\% power.

The dominant predictor of bruising severity was procedure code: procedure code~3 was associated with substantially higher depth ($\beta = 0.675$, $p = 0.006$) and extent scores across all models. This finding should inform stratification in future trial designs.

Forty-one per cent of the variance in depth scores was attributable to between-patient heterogeneity. The intended paired structure of the trial -- comparing the two eyes within each patient -- could not be recovered because treatment was assigned at the patient level rather than the eye level. Recovery of this within-patient contrast would have substantially increased statistical power and is the most important design change for any replication study.

Inter-rater reliability was acceptable for extent (ICC $= 0.888$) but below threshold for depth (ICC $= 0.515$), with a systematic upward bias from Grader~1. Future trials should include grader calibration sessions and additional raters.

%The study is still blinded at the time of this analysis. All treatment effects are reported as Condition~0 versus Condition~1. Unblinding should be performed only after all model specifications are locked, as documented in the pre-registered analysis plan.

In summary, the current evidence does not support a clinically important effect of the intervention on periorbital bruising, but the evidence base is insufficient to rule out a moderate or small effect. A properly powered, eye-level randomized trial with prospective recording of the quadrant-to-eye mapping is needed before any definitive clinical conclusion can be drawn.
